\chapter{Introdu\c{c}\~ao}
\label{cap:introducao}

\section{Motiva\c{c}\~ao}

O kernel Linux \'e um sistema operacional completo, com milh\~oes de linhas de
c\'odigo e d\'ecadas de evolu\c{c}\~ao. Estudar seu funcionamento a partir do
c\'odigo real exige percorrer camadas de abstra\c{c}\~ao, configura\c{c}\~oes de
compila\c{c}\~ao, subsistemas e caminhos de c\'odigo que n\~ao s\~ao essenciais
para entender o n\'ucleo do problema ``como um computador sai do boot e abre um
terminal?''.

O \textbf{Rust VGA Kernel} foi criado como um recorte did\'atico desse
problema. Ele reimplementa, em Rust, o caminho m\'inimo necess\'ario para:

\begin{enumerate}
  \item ser carregado por um bootloader padr\~ao (GRUB via Multiboot2);
  \item sair do modo de 32 bits e entrar em modo longo x86\_64;
  \item configurar p\'aginas e uma pilha de kernel;
  \item inicializar um console de texto VGA 80$\times$25;
  \item ler teclas do controlador PS/2 8042 e da serial COM1;
  \item instalar uma IDT, receber exce\c{c}\~oes e o timer do PIT;
  \item alocar frames de 4 KiB e materializar p\'aginas sob demanda;
  \item executar c\'odigo em ring 3 com separa\c{c}\~ao usu\'ario/kernel;
  \item oferecer um pequeno shell com comandos b\'asicos.
\end{enumerate}

O resultado n\~ao \'e um Linux reduzido, nem um porte linha a linha. \'E um
kernel independente, sem depend\^encias externas de Rust, que reproduz
\emph{conceitos} dos subsistemas equivalentes do Linux em uma escala
compreens\'ivel.

\section{Objetivo do documento}

Esta documenta\c{c}\~ao tem quatro objetivos:

\begin{itemize}
  \item \textbf{Explicar a arquitetura} do kernel: camadas, fluxo de boot,
        modelo de mem\'oria, modelo de execu\c{c}\~ao e depend\^encias entre
        m\'odulos.
  \item \textbf{Documentar as principais fun\c{c}\~oes} de cada m\'odulo,
        incluindo contratos, pr\'e-condi\c{c}\~oes, efeitos colaterais e
        limita\c{c}\~oes. O Ap\^endice~\ref{ap:funcoes} traz o cat\'alogo
        completo das 96 fun\c{c}\~oes do projeto.
  \item \textbf{Registrar o processo de build, execu\c{c}\~ao e valida\c{c}\~ao},
        de forma que outra pessoa consiga reproduzir o ambiente e os testes.
  \item \textbf{Comparar o projeto com o kernel Linux}, apresentando
        correspond\^encias de arquivos e fun\c{c}\~oes e uma estimativa
        transparente do n\'ivel de similaridade.
\end{itemize}

\section{Escopo}

\subsection*{O que est\'a inclu\'ido}

\begin{itemize}
  \item Boot BIOS por CD/ISO usando GRUB e o protocolo Multiboot2.
  \item Transi\c{c}\~ao de 32 bits para modo longo (x86\_64).
  \item GDT com segmentos ring 0/ring 3, TSS e pilhas de privil\'egio;
        tabelas de p\'agina est\'aticas no boot e metade alta do kernel.
  \item IDT com 256 portas, exce\c{c}\~oes, PIC remapeado e timer PIT em IRQ0.
  \item Alocador de frames f\'isicos de 4 KiB e \emph{demand paging} para
        regi\~oes de usu\'ario e de kernel.
  \item Demos em ring 3 com prote\c{c}\~ao U/S e \emph{syscall} did\'atica
        \code{int \$0x80}.
  \item Console VGA em modo texto, com cursor de hardware, rolagem e
        atributos de cor.
  \item Driver de teclado PS/2 por \emph{polling}, com Shift, Enter, Tab e
        Backspace.
  \item Driver da UART 16550 em COM1, usado para depura\c{c}\~ao e entrada.
  \item Shell sem aloca\c{c}\~ao din\^amica com dez comandos.
  \item \emph{Panic handler} que imprime a falha e para a CPU.
  \item Build reprodut\'ivel com Cargo, \code{ld.lld}, \code{grub-mkrescue}
        e QEMU.
\end{itemize}

\subsection*{O que n\~ao est\'a inclu\'ido}

\begin{itemize}
  \item Heap, alocador de objetos, \code{alloc} ou cole\c{c}\~oes din\^amicas.
  \item Processos completos, threads, escalonador, troca de contexto,
        espa\c{c}os de endere\c{c}o por processo ou um conjunto amplo de
        chamadas de sistema.
  \item Sistema de arquivos, rede, USB, ACPI, SMP ou suporte UEFI nativo.
  \item Drivers para outros v\'ideos al\'em do VGA texto legado.
  \item Endurecimento de seguran\c{c}a como SMEP, SMAP, NX, KASLR ou
        valida\c{c}\~ao de ponteiros de usu\'ario.
\end{itemize}

\section{P\'ublico-alvo}

O documento \'e destinado a:

\begin{itemize}
  \item estudantes de sistemas operacionais e arquitetura de computadores;
  \item desenvolvedores Rust interessados em \code{no\_std} e bare metal;
  \item pessoas que desejam entender o caminho entre BIOS, bootloader e
        console antes de estudar o kernel Linux em profundidade;
  \item avaliadores que precisem julgar o grau de correspond\^encia entre
        este projeto e o Linux.
\end{itemize}

\section{Conven\c{c}\~oes}

\begin{itemize}
  \item Nomes de c\'odigo aparecem em \code{fonte monoespa\c{c}ada}.
  \item Endere\c{c}os hexadecimais usam o prefixo \code{0x}.
  \item ``Metade alta'' refere-se a endere\c{c}os acima de
        \code{0xFFFF800000000000}, conven\c{c}\~ao usada pelo Linux e pelo
        alvo \code{x86\_64-unknown-none}.
  \item ``Original'' ou ``Linux'' refere-se, salvo indica\c{c}\~ao contr\'aria,
        \`a \'arvore Linux 7.3.1 usada como refer\^encia de compara\c{c}\~ao.
  \item Valores de endere\c{c}o e tamanho de se\c{c}\~ao citados como
        ``observados'' foram extra\'idos da compila\c{c}\~ao de refer\^encia e
        podem variar ligeiramente entre vers\~oes do compilador.
\end{itemize}

\section{Estado do projeto}

A vers\~ao documentada \'e a \code{0.1.0}, com as seguintes caracter\'isticas:

\begin{table}[H]
\centering
\caption{Resumo do estado da vers\~ao documentada.}
\label{tab:estado}
\begin{tabularx}{\linewidth}{@{}lX@{}}
\toprule
\textbf{Item} & \textbf{Situa\c{c}\~ao} \\
\midrule
Compila\c{c}\~ao & \texttt{make} conclui sem erros; alvo \code{x86\_64-unknown-none}. \\
Imagem & \code{vga-kernel.iso} El Torito, inicializ\'avel por BIOS. \\
Boot & Validado em QEMU com SeaBIOS e GRUB Multiboot2. \\
Console & VGA texto 80$\times$25, cursor, rolagem e cores. \\
Entrada & PS/2 por polling e serial COM1. \\
Interrup\c{c}\~oes & IDT com 256 portas, exce\c{c}\~oes, PIC remapeado e PIT a
~100 Hz. \\
Mem\'oria & Alocador de frames de 4 KiB, tabelas de 4 KiB e demand paging. \\
Prote\c{c}\~ao & GDT/TSS, ring 3 e bits U/S para usu\'ario e kernel. \\
Shell & \code{help}, \code{clear}, \code{echo}, \code{info}, \code{demand},
\code{usertest}, \code{userprot}, \code{mem}, \code{reboot}, \code{halt}. \\
Testes & Boot, comandos, entrada PS/2/serial, timer, page faults e ring 3
verificados. \\
Cobertura Linux & Subconjunto pequeno; similaridade baixa (Cap\'itulo~\ref{cap:similaridade}). \\
\bottomrule
\end{tabularx}
\end{table}
